\section{Discussion}
Group-level association and individual prediction answer different questions. Attendance, outside commitments, and desired-department match vary modestly across CGPA groups, but the WEKA questionnaire-only models remain close to the majority baseline. The best questionnaire-only 10-fold result is 30.04\% compared with 29.04\% for ZeroR. Students in different CGPA groups often provide similar questionnaire answers, so small distribution shifts do not create clean four-class boundaries.

Recent SGPA carries a more direct summary of academic history. When it is added to the main CGPA inputs, Random Forest reaches 43.92\% in 10-fold cross-validation. The readable J48 tree places recent SGPA at the root and uses outside commitments and desired-department match only within the highest recent-SGPA branch. That tree reaches 49.77\% on the one fixed test split. The cross-validation and fixed-split values use different model settings and should not be ranked as if they were the same experiment.

The fixed split also gives the questionnaire-only tree a small positive change from baseline, while the 10-fold comparison shows only a one-point best-model change. This variation reinforces the need to treat individual branches and single-split scores as exploratory. The stored confusion matrices, models, tree graphs, and raw WEKA logs make each reported value auditable.

The university imbalance limits how far the findings can generalise. The one-time self-reported design prevents causal conclusions. Questionnaire imputation was fitted before evaluation in the historical cleaning pipeline, so limited preprocessing leakage may make the scores optimistic. The results support language such as ``associated with,'' ``more common among,'' and ``helps prediction.'' They do not support claims that a questionnaire answer causes a grade or guarantees a student's outcome.

Future data collection should prioritise verified grades, course-level results, repeated observations of the same students, objective activity indicators, and better-balanced university sampling. A future evaluation should fit every imputation step inside each training fold before using the system for any decision beyond classroom demonstration.
